\section{Introduction}\label{sec:intro}

Autoregressive transformers cache the key and value projections of every
past token at every layer. This key--value (KV) cache makes decoding
linear rather than quadratic in context length, but it also grows without
bound: in half precision, a 32-layer 7B-class model spends roughly
$0.5$\,GB of accelerator memory per $1{,}000$ tokens, and the cache is
invariably the first resource to exhaust in long-context
serving~\cite{pope2023efficiently,kwon2023paged}. Because attention is
strongly sparse in practice, a long line of work evicts entries that are
unlikely to be attended again, retaining only a small fraction of the
cache under a \emph{memory contract}: if $\rho$ is the retention ratio,
the policy must keep at most $\rho L t$ entries in total across $L$ layers
at decoding step $t$~\cite{zhang2023h2o,liu2023scissorhands,%
xiao2024streamingllm,oren2024tova}.

Two design axes dominate the literature. \emph{Scoring} decides which
individual entries matter: cumulative attention mass
\cite{zhang2023h2o}, persistence of importance \cite{liu2023scissorhands},
pure recency \cite{xiao2024streamingllm}, or query-aware estimates
\cite{tang2024quest,li2024snapkv}. \emph{Allocation} decides how the
global budget is split across layers: uniform by default,
statically pyramidal in PyramidKV \cite{cai2024pyramidkv}, or dynamically
adapted to per-layer attention behaviour in CAKE \cite{qin2025cake} and
LAVa \cite{shen2025lava}. The allocation axis is attractive because
attention dispersion differs sharply across depth, and several studies
report sizeable quality gains from non-uniform budgets at equal total
memory.

\paragraph{The budget-consistency gap.}
This paper starts from an uncomfortable observation about the allocation
axis: \emph{none of the proposed layer-budget rules provably satisfies the
global memory contract that motivates them}. Static profiles are fixed
shapes that ignore both the contract and the runtime state; dynamic
heuristics cascade or re-balance budgets in ways that are evaluated
empirically but not guaranteed. The gap is not hypothetical. In an earlier
version of this work we ourselves proposed an entropy-exponent rule of the
form
$b_l(t)=\operatorname{clip}\!\bigl(\rho t\,(\kappa_l/\bar\kappa)^{\tau},
\,b_{\min}t,\,t\bigr)$,
where $\kappa_l$ measures attention dispersion in layer $l$ and
$\bar\kappa$ is its cross-layer mean. The rule looks innocuous: the
ratios $\kappa_l/\bar\kappa$ average to one. But the mean of a convex
power is not the power of the mean, and for the default $\tau=1.5$ the
Jensen gap is systematic. In simulation over drifting dispersion signals,
the formula \emph{as written} violates the advertised contract at
$97.4\%$ of steps, with a mean overshoot of $+4.2\%$ and worst-case
$+10.6\%$ (Figure~\ref{fig:budget}); at $20\%$ nominal retention this
silently reclaims a nontrivial fraction of the savings that justify the
method. An implementer is then forced to choose between two undocumented
repairs---renormalising after exponentiation, or clipping and hoping---and
the published formula describes neither. Similar exponent-style rules
appear across the adaptive-allocation literature, usually without a
statement of which repair, if any, is applied.

\paragraph{This work.}
We treat budget consistency as a correctness property, not an
implementation detail, and rebuild the eviction stack around it.
\sage{} keeps the three mechanisms of the original design---decayed
cumulative scoring (\acs{}), layer-adaptive budgets driven by measured
dispersion (\laba{}), and protection of sinks and recency (\dre{})---but
replaces every heuristic allocation step with a single primitive: the
projection of a weight-proportional ideal allocation onto the capped
simplex $\{b\in\mathbb{Z}^L:\ \ell_l\le b_l\le u_l,\ \sum_l b_l=B\}$ by
bisection on a water level. We prove that this \emph{capped water-filling}
projection hits the global budget exactly at every step, respects
per-layer floors that encode protected regions, and costs
$O(L\log B)$ (Theorem~\ref{thm:exact}). The same primitive is used for
\emph{every} method in our comparison, so StreamingLLM, H2O, and
PyramidKV-style baselines all hold identical retained-entry counts at
every step of every run---a matched-entry protocol that removes the
nominal-versus-actual confound we criticise in prior evaluations,
including our own earlier one.

\paragraph{Contributions.}
Our contributions are as follows.
\begin{enumerate}[leftmargin=1.6em,itemsep=2pt,topsep=3pt]
\item \textbf{Diagnosis.} We quantify the Jensen overshoot of
entropy-exponent budget formulas ($97.4\%$ violating steps; $+4.2\%$ mean
excess in simulation), and identify budget consistency as a missing
correctness property across adaptive layer-allocation methods
(Section~\ref{sec:prelim}).
\item \textbf{Exact-budget allocation.} We give a capped water-filling
projection with a full correctness proof, integer-exact budgets,
protection floors, warmup analysis, and $O(L\log B)$ complexity
(Theorem~\ref{thm:exact}, Section~\ref{sec:method}).
\item \textbf{Analysis of scoring and weighting.} We prove an output
perturbation bound proportional to evicted attention mass
(Lemma~\ref{lem:perturb}), an explicit two-token separation showing when
decayed scoring strictly improves on cumulative scoring
(Proposition~\ref{prop:staleness}), and an equal-marginal argument that
motivates dispersion-weighted budgets under a stylised power-law tail
model (Proposition~\ref{prop:marginal}).
\item \textbf{Matched-entry, fully traceable evaluation.} All baselines
run through the same projector and cache manager on two attention
architectures---GPT-2 (multi-head) and Qwen2.5-0.5B-Instruct
(grouped-query)---with WikiText-2 perplexity, budget sweeps, needle
retrieval with per-layer retention diagnostics, a multi-instruction
probe, sensitivity analyses, and per-step accounting logs. Every number in
this paper is generated by the released code from the released logs; the
synthetic records of our earlier manuscript are discarded entirely
(Section~\ref{sec:exp}).
\item \textbf{Calibrated findings.} Dispersion-adaptive allocation beats
static-pyramid allocation on both architectures and beats uniform
allocation on GPT-2; sink protection is existential (its removal
multiplies perplexity by $84\times$); the decay rate trades language
modelling quality against retrieval robustness; and history-only scoring
loses mid-context needles at $20\%$ budget---a fundamental limitation we
measure directly through cache-level retention rather than only through
end accuracy (Section~\ref{sec:disc}).
\end{enumerate}

We deliberately report the study at the scale it was actually run---CPU
inference, models up to $0.5$B parameters, contexts up to $4{,}096$
tokens---and release a GPU replication protocol rather than extrapolate.
To our knowledge this is the first KV-eviction study whose memory-contract
claims are backed by per-step, machine-checked accounting of every
configuration reported.
